%-*-tex-*-
\ifundefined{writestatus} \input status \relax \fi %
\chcode{shead}
\def\cqu{\cquote{The chief virtue that language can have is clearness, and
nothing detracts from it so much as the use of unfamiliar
words.}{On the Natural Faculties, Galen }
}

\chapterhead{shead}{CHAPTER\cr and SECTION\cr HEADS}
\intex\ supplies commands for creating chapter and section heads. The default
style of the  heads are those given in this book. They form the basis
of the style that naturally results from the use of \intex. Some changes in
style such as the fonts used are quite easy; others require more care. 

This chapter introduces the commands for creating section heads, whether they
occur during the initial part of the document or within appendices. All of
the commands have the same format with or without autonumbering or automatic
generation of table of contents in force. 

\intex\ considers that a document consists of a {\it prelude} part that will
include the preface, table of contents, and other special summaries, an {\it
interlude}
consisting of chapters and sections, appendices, which may appear at the
end of a document or at the end of a chapter, and a {\it postlude} which
would include, perhaps, an index, references, and even appendices. The {\it
prelude}, appendices, and {\it postlude} are
enclosed by a |\begin...\end...| pair that form a
group and set up the styles appropriate for that section. 

The forms that introduce sections, tables of contents, or chapters all have
commands of the form |\...head|, such as |\shead|, 
where the |head| is supposed invoke images
of starting a unit. In \intex, all commands that have a |\begin...| have an
accompanying |\end...| and all create a group. The |\...head| commands are
key elements in controlling the setting and incrementing of the various
numbers used in the document. In addition, they have a relatively complex,
but accessible format part that can be modified if required. 

The section head commands  |\chead|, 
|\shead|, |\sshead|, |\ssshead|, and
|\dssshead| stand for chapterhead, 
sectionhead, subsectionhead, subsubsectionhead,
and diminished subsubsectionhead. This is the only example of commands in
\intex\ for which the abbreviations exist without the full forms. This is
primarily for typing and visual convenience in the original file. The fonts
used have the same abbreviated form. 

If an |\shead| is immediately followed by an |\sshead| or any other
sectionhead form, there is a chance that they may be split apart if they
occur at the bottom of the page. {\bf To prevent this, and to
make the spacing reasonable,  the command
|\nosheadbreak| must be placed in between the two commands.}
The general spacing parameters are given in Section \ref{sheadspac}.

If the |.toc| file is open for making a table of contents list, each section
head form writes information to this file. The actual form sent is determined
by the definition of a |\...tocout| form. For instance the |\shtocout| is
defined as |\def\shtocout{\string\shtoc}|. The |\shtoc| is an \intex\ form
that produces a section head line in the table of contents. 


\shead{sheadcomlist}{Command List}
\beginthreecolumn
\hfuzz=10pt
\ext|\Alphabetic|
\ext|\alphabetic|
\ext|\Appendix|
\ext|\beginappendices|
\ext|\beginprelude|
\ext|\Chapter|
\ext|\chapterstartform|
\ext|\chead|
\ext|\cheadfont|
\ext|\chnum|
\ext|\chtagreplaceformat|
\ext|\chnumform|
\ext|\chtocout|
\ext|\endappendices|
\ext|\endprelude|
\ext|\dssshead|
\ext|\dsssheadfont|
\ext|\dssshtocout|
\ext|\innerchead|
\ext|\innerdssshead|
\ext|\innershead|
\ext|\innersshead|
\ext|\innerssshead|
\ext|\nosheadbreak|
\ext|\preludehead|
\ext|\prtocout|
\ext|\shead|
\ext|\sheadfont|
\ext|\shnum|
\ext|\shnumform|
\ext|\shtagreplaceformat|
\ext|\shtocout|
\ext|\sshead|
\ext|\ssheadfont|
\ext|\sshnum|
\ext|\sshnumform|
\ext|\sshtocout|
\ext|\ssshead|
\ext|\sssheadfont|
\ext|\ssshnum|
\ext|\ssshnumform|
\ext|\ssshtocout|
\ext|\UpperRoman|
\endthreecolumn

\shead{sheadcomforms}{Command Forms --- Making Head Forms}
The commands in this section are very similar in form  and will be described
in together. There are six predefined |<head types>|: |prelude|, |c| 
(chapter),
|s|, |ss|, |sss|, and |dsss|. All but |prelude| and |dsss| take two arguments
in their |\<head type>head| form. 

Number format and font commands  are discussed in Section \ref{numforms}.

\beginblockmode
\mbr
\ext\@|\chead{<number or tag>}{<title>}|
\nbr
\ext\@|\preludehead{<title>}|
\nbr
\ext\@|\shead{<number or tag>}{<title>}|
\nbr
\ext\@|\sshead{<number or tag>}{<title>}|
\nbr
\ext\@|\ssshead{<number or tag>}{<title>}|
\nbr
\ext\@|\dssshead{<title>}|
\nbr
All of these commands create the various forms of section or chapter heads
indicated by the command name. The |<number or tag>| field is interpreted as a
tag if |autonumbering| is on or inserted as is into the appropriate place if
is not. The lack of this field in |\dssshead| and |\preludehead| implies that
it is impossible to tag or refer to these document parts except through
their, presumably known, titles. The tag replacement format  and associated 
|\<head type>hnumform| are 
discussed separately in
Section \ref{numforms}. 

There are some restrictions in the use of the |<title>| field. This should
{\bf not} include |\footnote|, |\cite|, or |\ref| commands. If it is
absolutely necessary to use these, see Chapter~{\ref{footnot}} to see how to
get around the |\footnote| restriction, and use |\quietref{<tag>}| 
for the other
two. {\it Quiet} is not quite {\it silent}.

\mbr
\ext\@|\Appendix|
\nbr
This command defines the actual form that will be seen in an |\shead| used for
appendices. It is defined as ``Appendix'' in |\englishversion| and
``Appendice'' in |\versionfrancaise|.

\mbr
\ext\@|\beginappendices  \endappendices|
\nbr
These commands are intended to enclose the appendices that occur at the end
of a chapter or book. They reset the {\it special} auto numbers, and change
the format of the numbers when an |\shead| or its relatives are called.
|\chead| should not be called in an appendix. 
Since the |\begin...\end...| pair form a group, 
changes that only affect the appendices are
easily made.

\mbr
\ext\@|\beginprelude  \endprelude|
\nbr
These commands are intended to enclose the prelude that occurs at the 
the start of a document. Preludes usually include such entities as prefaces,
tables of contents, figures, and tables, acknowledgments, distribution lists
or perhaps executive summaries. 
Since the |\begin...\end...| pair form a group, changes that only affect the
prelude parts are 
easily made.
\mbr
\ext\@|\Chapter|
\nbr
This command defines  the actual form used in a |\chead|.
It is defined as ``Chapter'' in |\englishversion| and
``Chapitre'' in |\versionfrancaise|.

\mbr
\ext\@|\chapterstartform = {<special pre chapter tests>}|
\nbr
The token list for this command is executed inside |\chead| before the
first page of the chapter is produced. For this book it was used to check
whether the chapter starting page was odd. If it was not, a ``quote'' page
was produced. 

\endblockmode

\shead{numforms}{Command Forms --- Number Formats}
The number formats when autonumbering is in force are defined using a few
commands. First the fonts in the various head forms are called
\@|\<head type>headfont|. These are defined when a document font size (see
Chapter \ref{fonts}) is defined. For example
|\documentfont{\tenpoint}| determined the document font family for this book.
The individual fonts may be changed, although it is not recommended, by
changing the |\<head type>headfont|. For example, 
|\let\sheadfont = \twelvedu| will change the section head font for the rest
of the document, or until the end of the group if it is called within a
group. A change of this sort will not modify the other fonts.

The number formats are modified by changing the |\<head type>hnumform| such as 
|\chnumform|, |\shnumform|, |\sshnumform|, |\ssshnumform|,
respectively for the chapter  and the various section heads.
\@|\Alphabetic{<integer>}|, \@|\alphabetic{<integer>}|,  and
\@|\UpperRoman{<integer>}| are used to change the format of an integer
(counter) to either the equivalent uppercase letter, lowercase letter
or uppercase roman numeral.  
The \intex\ forms for |\<head type>hnumform| are
\begintt
\def\<head type>hnumform{\the\<head type>hnum}
\endtt
for all |<head type>| except that in appendices |\shnumform| is defined to be
\begintt
\def\shnumform{\Alphabetic{\the\shnum}}
\endtt
This effectively changes the appendix numbering to be roman capitals.

The command \@|\chtagreplaceformat| is the actual command that is used to
place the ``Chapter \ref{shead}'' at the beginning of this chapter. The
two definitions involved in this replacement are
\begintt
\def\chtagreplaceformat {Chapter \chtagrefformat}
\def\chtagrefformat {\chnumform}
\endtt
The \@|\shtagreplaceformat| is
defined similarly and is used in |\sheads|. 

A change in the form of the command written to the table of contents file may
be obtained simply by redefining the appropriate |\...tocout|. 


\shead{sheadspac}{Command Forms --- Section Head Spacing}
The spacing both before and after the various |\...head|s are determined by
the following parameters. Stretch and shrink glue should be specified to give
\tex\ some freedom to break pages. 
\bshortcomlist
\ext\@|\prsheadskip [=] <glue>|& pre shead skip \cr 
\ext\@|\posheadskip [=] <glue>|& post shead skip \cr
\ext\@|\prssheadskip [=] <glue>|& pre sshead skip \cr
\ext\@|\possheadskip [=] <glue>|& post sshead skip \cr
\ext\@|\prsssheadskip [=] <glue>|& pre ssshead skip \cr
\ext\@|\posssheadskip [=] <glue>|& post ssshead skip \cr
\ext\@|\prdsssheadskip [=] <glue>|& pre dssshead skip \cr
\ext\@|\podsssheadskip [=] <glue>|& post dssshead skip \cr
\eshortcomlist


\shead{genform}{Major Modifications of the Section Head Formats}
\intex\ has separate |\cheadformat|, \@|\preludeheadformat|, and a very
general section head format. These are easily modifiable without upsetting
any of the houskeeping. Any format that satisfies the input parameter
constraints will work. 

\beginblockmode
\mbr
\ext\@|\cheadformat{<chapter>}{<text>}|
\nbr
The |<chapter>| is the chapter identifier. In this chapter it was 
{\bf Chapter~\ref{shead}}. The |<text>| is title of the chapter. The
particular form was built using |\vcenter| and |\halign|. 

\mbr
\ext\@|\gensheadformat{<font>}{<number>}{<text>}{<separator>}|
\nbr
This command uses four parameters. The first is the font, the second the
number, the third is the title of the section and the fourth is the separator
between the number and the text. The form that is used checks the number to
see if it is empty. If it is, the |<text>| is placed flush left. The |<text>|
is actually placed in a |\vbox| with |\veryraggedright| spacing so that there
will be no hyphenation and long titles will automatically split. Details are
in {\tt secthead.tex}. 

\mbr
\ext\@|\preludheadformat{<title>}|
\nbr
This command is defined with one parameter that is assumed to be the title of
the particular section. Since this one is quite simple the default definition
is given as 
\begintt
\def\preludeheadformat#1{\centerline{\sheadfont #1}\vskip 1cm}
\endtt
\endblockmode
\ejectpage
\done





